\documentclass{article}
\usepackage{graphicx}
\usepackage{amsmath}
\usepackage{amssymb}
\usepackage{amsfonts}
\usepackage{hyperref}
\usepackage{exercises}
\usepackage{indentfirst}
\usepackage{circledsteps}
\usepackage{amsthm}
\usepackage{csquotes}
\usepackage{yhmath}
\usepackage{comment}
\usepackage{placeins}
\usepackage{geometry}
\usepackage{hyperref}
\usepackage{tikz}
\usepackage{tabto}
\newcommand*{\oldneg}{\mathord{\sim}}
\newcommand{\Z}{\mathbb{Z}}
\newcommand{\N}{\mathbb{N}}
\let\oldemptyset\emptyset

\theoremstyle{definition}
\newtheorem{definition}{Definition}
\newtheorem{theorem}{Theorem}[section]
\newtheorem{corollary}{Corollary}[theorem]
\newtheorem{lemma}[theorem]{Lemma}

\setlength{\parindent}{.5in}
\geometry{margin=0.5in}
\setlength{\parskip}{12pt}
\title{Chapter 12 Notes}
\begin{document}
\date{}
\maketitle

\section*{Section 0}

\subsection*{Notes}
Note that when the relation $x=y$ is used (for instnace in $P_1$), this is a substitute for the formula
\[ \forall z(z\in x\equiv z\in y) \]

\section*{Section 1}

\subsection*{Exercise 1.1}
We show that for finite $\alpha$, $L_\alpha=R_\alpha$.
\begin{proof}
By induction. $L_0=R_0$ by definition. Suppose $L_n=R_n$ where $n$ is finite. By $C_0$, $L_n$ is transitive
(and thus also extensional), and hence by $P_3$, $\mathcal{F}(L_n)=\mathcal{P}(L_n)=\mathcal{P}(R_n)=R_{n+1}$, 
therefore completing the induction.
\end{proof}

\subsection*{Exercise 1.2}
By Exercise 1.1, for finite $\alpha$, $L_\alpha=R_\alpha$. Then it obviously follows that 
$L_\omega=R_\omega$. However, $L_{\omega+1}\ne R_{\omega+1}$. For recall that, if 
a set $A$ is infinite, then $\mathcal{F}(A)\preceq A$. 
Thus $L_{\omega+1}=\mathcal{F}(L_\omega)\preceq L_\omega$ (because $L_\omega$ is infinite). But
$R_{\omega+1}=\mathcal{P}(R_\omega)\succ R_\omega$ ($\succ L_\omega$) by 
Cantor's Theorem. Hence $L_{\omega+1}\ne R_{\omega+1}$, and $R_{\omega+1}\succ L_{\omega+1}$.
\end{document}
